\clearpage
\section[Анализ предметной области и постановка задачи]{\centering\MakeUppercase{Анализ предметной области и постановка задачи}}

\subsection{\centering{Предметная область}}

Паллетизация в складской логистике состоит в формировании устойчивой грузовой единицы из множества товарных коробов. На практике решение должно быть не только плотным, но и выполнимым для транспортировки: тяжелые товары не должны разрушать хрупкие, нестабильные слои не должны приводить к падению груза, а ограничения по высоте и массе должны соблюдаться для совместимости со складским оборудованием и транспортом.

В условии Smart 3D Packing отдельно выделены типичные проблемы ручной или простой эвристической сборки: низкая утилизация объема, риск повреждений из-за неправильного распределения веса и нестабильность конструкции. Эти факторы показывают, что задача имеет многокритериальный характер. Один только максимум занятого объема не гарантирует хорошего решения, если оно нарушает устойчивость, портит товар или требует слишком долгого расчета.

С математической точки зрения задача относится к семейству bin packing и cutting/packing problems~\cite{lodi2002}. В классической трехмерной постановке требуется разместить прямоугольные объекты в одном или нескольких контейнерах. Даже более простые варианты bin packing являются вычислительно сложными, а добавление трехмерной геометрии и физических ограничений резко увеличивает размер пространства решений~\cite{martello2000,lodi2002}.

\subsection{\centering{Входные данные}}

Вход одной задачи представлен JSON-объектом. Он содержит описание паллеты и список типов коробов. Для паллеты задаются:
\begin{itemize}
  \item длина \code{length\_mm};
  \item ширина \code{width\_mm};
  \item максимальная высота укладки \code{max\_height\_mm};
  \item максимальная масса \code{max\_weight\_kg}.
\end{itemize}

Для каждого типа короба задаются идентификатор SKU, описание, габариты, масса, количество экземпляров и дополнительные признаки:
\begin{itemize}
  \item \code{strict\_upright} -- запрет поворота вокруг горизонтальных осей;
  \item \code{fragile} -- признак хрупкости;
  \item \code{stackable} -- возможность ставить другие короба поверх данного.
\end{itemize}

В генераторе проекта используются типовые архетипы товаров фуд-ритейла: вода, сахар, вино, бананы, чипсы, яйца и консервы. Размеры и массы этих архетипов варьируются небольшим случайным шумом, что позволяет получать воспроизводимые, но не полностью одинаковые задачи.

\subsection{\centering{Выходные данные}}

Результатом работы решателя является JSON-ответ, содержащий идентификатор задачи, версию решателя, время расчета и список размещенных коробов. Для каждого размещенного экземпляра указываются:
\begin{itemize}
  \item \code{sku\_id} и \code{instance\_index};
  \item координаты нижнего ближнего левого угла \code{x\_mm}, \code{y\_mm}, \code{z\_mm};
  \item фактические размеры после поворота;
  \item код ориентации \code{rotation\_code}.
\end{itemize}

Если часть объектов не размещена, она фиксируется в массиве \code{unplaced} с указанием количества и причины: превышение массы, высоты или отсутствие места. Такая структура удобна для последующей валидации и анализа: ответ содержит как геометрию решения, так и информацию о неполноте размещения заказа.

\subsection{\centering{Жесткие ограничения}}

В проекте реализован валидатор \code{validator.py}, который проверяет физическую корректность решения. Нарушение любого жесткого ограничения делает решение невалидным для конкретной задачи.

Основные ограничения следующие:
\begin{enumerate}
  \item каждый размещенный короб должен целиком находиться внутри габаритов паллеты;
  \item объемы любых двух коробов не должны пересекаться;
  \item суммарная масса размещенных товаров не должна превышать грузоподъемность паллеты;
  \item если \code{strict\_upright=true}, фактическая высота короба должна совпадать с исходной высотой;
  \item короб, расположенный выше пола паллеты, должен иметь опору снизу не менее чем на 60\% площади основания согласно формуле~\eqref{eq:support};
  \item фактические размеры размещенного короба должны быть перестановкой исходных размеров, что исключает некорректное изменение габаритов.
\end{enumerate}

Условие опоры записывается следующим образом:
\begin{equation}
\frac{S_{\mathrm{support}}}{S_{\mathrm{base}}} \geq 0.6,
\label{eq:support}
\end{equation}
где $S_{\mathrm{base}}$ -- площадь основания проверяемого короба, а $S_{\mathrm{support}}$ -- суммарная площадь пересечения его основания с верхними гранями коробов, лежащих непосредственно ниже.

\subsection{\centering{Метрика качества}}

Если решение прошло жесткие проверки, валидатор вычисляет итоговую метрику:
\begin{equation}
\mathrm{score}=0.50U_{\mathrm{vol}}+0.30C_{\mathrm{items}}+0.10F_{\mathrm{frag}}+0.10T_{\mathrm{time}}.
\label{eq:score}
\end{equation}

Компонента утилизации объема равна
\begin{equation}
U_{\mathrm{vol}}=\frac{\sum_{i \in P} l_i w_i h_i}{LWH},
\label{eq:volume}
\end{equation}
где $P$ -- множество размещенных коробов, $l_i,w_i,h_i$ -- их фактические размеры, $L,W,H$ -- габариты паллеты.

Полнота заказа определяется как
\begin{equation}
C_{\mathrm{items}}=\frac{N_{\mathrm{placed}}}{N_{\mathrm{total}}},
\label{eq:coverage}
\end{equation}
где $N_{\mathrm{placed}}$ -- число размещенных экземпляров, $N_{\mathrm{total}}$ -- общее число экземпляров во входной задаче.

Оценка хрупкости задается формулой
\begin{equation}
F_{\mathrm{frag}}=\max(0, 1-0.05n_{\mathrm{viol}}),
\label{eq:frag}
\end{equation}
где $n_{\mathrm{viol}}$ -- число случаев, когда короб массой более 2 кг расположен непосредственно над хрупким коробом с ненулевым перекрытием в плоскости $XY$.

Временной балл имеет ступенчатый вид: при времени до 1 секунды он равен 1.0, при времени от 1 до 5 секунд -- 0.7, при времени от 5 до 30 секунд -- 0.3, при большем времени -- 0.0. Такая метрика поощряет алгоритмы, которые дают хороший результат в режиме, близком к реальному времени.

\subsection{\centering{Критерии качества упаковки}}

Формула~\eqref{eq:score} показывает, что 80\% итогового score формируют плотность упаковки и доля размещенных товаров. Однако оставшиеся 20\% также существенны, поскольку они отражают практическую применимость решения. Метод, который достигает высокой плотности за счет нарушения хрупкости или слишком большого времени расчета, может уступить более сбалансированному подходу.

Таким образом, в работе используются следующие критерии сравнения:
\begin{itemize}
  \item средний итоговый score на наборе задач;
  \item средняя утилизация объема;
  \item средняя полнота размещения;
  \item средний балл хрупкости;
  \item средний временной балл и фактическое время решения;
  \item число валидных решений.
\end{itemize}

\subsection{\centering{Выводы по разделу}}

Задача Smart 3D Packing является прикладной модификацией 3D Bin Packing, в которой геометрическая упаковка дополняется логистическими ограничениями. Наиболее важными особенностями постановки являются поддержка поворотов, проверка опоры, учет массы, запрет некорректного размещения тяжелых товаров над хрупкими и ограничение времени. Это делает задачу удобной для сравнения как быстрых эвристик, так и более сложных метаэвристических и ML-подходов.
